\documentclass[12pt,a4paper]{article}

\usepackage[utf8]{inputenc}
\usepackage[T1]{fontenc}
\usepackage[margin=2.5cm]{geometry}
\usepackage{hyperref}
\usepackage{xcolor}
\usepackage{listings}
\usepackage{booktabs}
\usepackage{tabularx}
\usepackage{array}
\usepackage{enumitem}
\usepackage{titlesec}
\usepackage{parskip}
\usepackage{fancyhdr}
\usepackage{lmodern}
\usepackage{microtype}
\usepackage{mdframed}

% Cores
\definecolor{codebg}{HTML}{F4F4F4}
\definecolor{codeframe}{HTML}{CCCCCC}
\definecolor{linkcolor}{HTML}{1A56DB}
\definecolor{notecolor}{HTML}{FFF8E1}
\definecolor{noteframe}{HTML}{F59E0B}
\definecolor{sectioncolor}{HTML}{1E3A5F}

% Hyperlinks
\hypersetup{
  colorlinks=true,
  linkcolor=linkcolor,
  urlcolor=linkcolor,
  pdftitle={Guia de Alterações — BH News Chatbot},
  pdfauthor={BH News Chatbot},
}

% Código
\lstdefinestyle{bash}{
  backgroundcolor=\color{codebg},
  frame=single,
  rulecolor=\color{codeframe},
  basicstyle=\ttfamily\small,
  breaklines=true,
  breakatwhitespace=false,
  keepspaces=true,
  columns=flexible,
  showstringspaces=false,
  commentstyle=\color{gray}\itshape,
  morecomment=[l]{\#},
  xleftmargin=8pt,
  xrightmargin=8pt,
  aboveskip=6pt,
  belowskip=6pt,
}
\lstset{style=bash}

% Títulos
\titleformat{\section}{\large\bfseries\color{sectioncolor}}{}{0em}{}[\titlerule]
\titleformat{\subsection}{\normalsize\bfseries}{}{0em}{}
\titleformat{\subsubsection}{\normalsize\itshape}{}{0em}{}
\titlespacing{\section}{0pt}{18pt}{8pt}
\titlespacing{\subsection}{0pt}{12pt}{4pt}

% Cabeçalho/rodapé
\pagestyle{fancy}
\fancyhf{}
\fancyhead[L]{\small\textit{BH News Chatbot}}
\fancyhead[R]{\small\textit{Guia de Alterações}}
\fancyfoot[C]{\small\thepage}
\renewcommand{\headrulewidth}{0.4pt}

% Nota de destaque
\newenvironment{nota}{%
  \begin{mdframed}[backgroundcolor=notecolor,linecolor=noteframe,linewidth=1.5pt,
    leftmargin=0pt,rightmargin=0pt,innerleftmargin=10pt,innerrightmargin=10pt,
    innertopmargin=8pt,innerbottommargin=8pt]
  \small
}{%
  \end{mdframed}
}

% Tabela com fundo cinza no header
\newcolumntype{L}{>{\raggedright\arraybackslash}X}

\begin{document}

% ----- Capa -----
\begin{titlepage}
  \centering
  \vspace*{3cm}
  {\Huge\bfseries\color{sectioncolor} Guia de Alterações\par}
  \vspace{0.5cm}
  {\Large BH News Chatbot\par}
  \vspace{1.5cm}
  \rule{0.5\linewidth}{0.4pt}\par
  \vspace{1cm}
  {\large Este documento descreve as funcionalidades adicionadas recentemente ao projeto,
  ensina como usá-las e responde às dúvidas mais comuns.\par}
  \vfill
  {\small Repositório: \href{https://github.com/paraenseembh/opensource_bh}%
    {github.com/paraenseembh/opensource\_bh}\par}
  \vspace{0.3cm}
  {\small Branch: \texttt{claude/bh-news-chatbot-NWtrw}\par}
  \vspace{1cm}
\end{titlepage}

% ----- Sumário -----
\tableofcontents
\newpage

% ===========================================================================
\section{O que foi adicionado}

\begin{tabularx}{\linewidth}{lLl}
  \toprule
  \textbf{Alteração} & \textbf{Descrição} & \textbf{Arquivo} \\
  \midrule
  Script de verificação & Testa conectividade, \texttt{complete()}, \texttt{stream()} e
    tratamento de erros de cada provedor & \texttt{test\_api.py} \\
  \addlinespace
  Testes de erros & Verifica que o sistema rejeita chaves inválidas, modelos inexistentes
    e provedores desconhecidos & \texttt{test\_api.py} \\
  \addlinespace
  Maritaca (Sabiá) & Terceiro provedor de LLM, focado em português brasileiro
    & \texttt{llm\_provider.py} \\
  \addlinespace
  Dependência \texttt{openai} & Necessária para comunicar com a API compatível da Maritaca
    & \texttt{requirements.txt} \\
  \addlinespace
  Flag \texttt{-{}-maritaca-key} & Permite passar a chave da Maritaca pela linha de comando
    & \texttt{main.py} \\
  \bottomrule
\end{tabularx}

% ===========================================================================
\section{Script de verificação de APIs (\texttt{test\_api.py})}

\subsection{Para que serve}

Antes de rodar o chatbot, use este script para confirmar que suas chaves de API estão
corretas e que os provedores estão respondendo. Ele executa dois grupos de testes:

\begin{itemize}
  \item \textbf{Testes funcionais:} chamam a API de verdade com uma mensagem simples.
    Precisam de chave válida.
  \item \textbf{Testes de erros esperados:} verificam que o sistema rejeita situações
    inválidas corretamente. \textbf{Não precisam de chave.}
\end{itemize}

\subsection{Como executar}

\begin{lstlisting}
# Testa todos os provedores configurados
python bh_news_chatbot/test_api.py

# Testa apenas um provedor
python bh_news_chatbot/test_api.py --provider anthropic
python bh_news_chatbot/test_api.py --provider gemini
python bh_news_chatbot/test_api.py --provider maritaca

# Roda apenas os testes de erros (nao precisa de nenhuma chave)
python bh_news_chatbot/test_api.py --only-errors
\end{lstlisting}

\subsection{Como interpretar a saída}

\begin{lstlisting}
[check]  verde    -- teste passou
[x]      vermelho -- teste falhou (veja o motivo na mesma linha)
[-]      amarelo  -- teste pulado (motivo entre parenteses)
\end{lstlisting}

\textbf{Exemplo de saída com chave configurada:}

\begin{lstlisting}
Maritaca
-------------------------------------------------
  testes funcionais
  [ok] create_provider()   Maritaca (sabia-4)
  [ok] complete()          1.3s  "OK"
  [ok] stream()            1.1s  4 token(s)  "OK"

  erros esperados
  [ok] chave vazia -> ValueError
  [ok] chave invalida -> complete() falha
  [ok] chave invalida -> stream() falha
  [ok] modelo inexistente -> complete() falha
\end{lstlisting}

\textbf{Exemplo sem chave configurada:}

\begin{lstlisting}
Maritaca
-------------------------------------------------
  testes funcionais
  [-]  todos os testes funcionais (MARITACA_KEY nao configurada)

  erros esperados
  [ok] chave vazia -> ValueError
  [ok] chave invalida -> complete() falha  (ModuleNotFoundError)
  [ok] chave invalida -> stream() falha
  [-]  modelo inexistente -> complete() falha (requer chave valida)
\end{lstlisting}

\begin{nota}
Os testes funcionais ficam com \texttt{-{}-} se a chave não estiver configurada. Isso é
normal --- configure a chave e rode novamente.
\end{nota}

\subsection{O que cada teste verifica}

\begin{tabularx}{\linewidth}{lL}
  \toprule
  \textbf{Teste} & \textbf{O que valida} \\
  \midrule
  \texttt{create\_provider()} & A fábrica cria o objeto do provedor sem erros \\
  \texttt{complete()} & A API responde a uma mensagem simples \\
  \texttt{stream()} & A API retorna tokens em streaming \\
  chave vazia $\to$ ValueError & O sistema rejeita chave em branco antes de chamar a API \\
  chave inválida $\to$ complete() falha & A API retorna erro de autenticação para chave falsa \\
  chave inválida $\to$ stream() falha & Idem, no modo streaming \\
  modelo inexistente $\to$ complete() falha & A API retorna erro para nome de modelo inválido \\
  \bottomrule
\end{tabularx}

% ===========================================================================
\section{Integração com Maritaca.ai (Sabiá)}

\subsection{O que é a Maritaca}

A \href{https://maritaca.ai}{Maritaca AI} é uma empresa brasileira que desenvolve modelos
de linguagem otimizados para o português. Os modelos da família \textbf{Sabiá} têm
desempenho superior em textos em português em relação a modelos genéricos do mesmo tamanho.

\subsection{Modelos disponíveis}

\begin{tabularx}{\linewidth}{llL}
  \toprule
  \textbf{Modelo} & \textbf{Custo} & \textbf{Uso recomendado} \\
  \midrule
  \texttt{sabia-4} & Mais alto & Chat principal, respostas de maior qualidade \\
  \texttt{sabiazinho-4} & Mais baixo & Categorização, tarefas rápidas, menor latência \\
  \bottomrule
\end{tabularx}

\subsection{Como obter a chave}

\begin{enumerate}
  \item Acesse \href{https://plataforma.maritaca.ai/chaves-de-api}{plataforma.maritaca.ai/chaves-de-api}
  \item Crie uma conta (novos usuários recebem R\$\,20 de crédito)
  \item Gere uma nova chave de API
\end{enumerate}

\subsection{Como configurar}

\begin{lstlisting}
# Variavel de ambiente (recomendado)
export MARITACA_KEY=sua-chave-aqui

# Ou passe direto na linha de comando
python bh_news_chatbot/main.py --provider maritaca --maritaca-key sua-chave-aqui
\end{lstlisting}

\subsection{Como usar no chatbot}

\begin{lstlisting}
# Uso basico com Sabia-4
python bh_news_chatbot/main.py --provider maritaca

# Com modelo rapido (menor custo)
python bh_news_chatbot/main.py --provider maritaca --model sabiazinho-4

# Com categorizacao e 5 artigos por area
python bh_news_chatbot/main.py --provider maritaca --max-per-area 5 --categorize

# Apenas categorizar sem abrir o chat
python bh_news_chatbot/main.py --provider maritaca --only-categorize
\end{lstlisting}

\subsection{Verificar se a chave está funcionando}

\begin{lstlisting}
export MARITACA_KEY=sua-chave-aqui
python bh_news_chatbot/test_api.py --provider maritaca
\end{lstlisting}

% ===========================================================================
\section{Comparativo de provedores}

\begin{tabularx}{\linewidth}{lllLL}
  \toprule
  \textbf{Provedor} & \textbf{Env var} & \textbf{Modelo chat} & \textbf{Modelo rápido} & \textbf{Foco} \\
  \midrule
  Anthropic & \texttt{ANTHROPIC\_API\_KEY} & \texttt{claude-sonnet-4-6} & \texttt{claude-haiku-4-5} & Geral \\
  Google Gemini & \texttt{GEMINI\_API\_KEY} & \texttt{gemini-2.0-flash} & \texttt{gemini-2.0-flash} & Geral, multimodal \\
  Maritaca & \texttt{MARITACA\_KEY} & \texttt{sabia-4} & \texttt{sabiazinho-4} & Português BR \\
  \bottomrule
\end{tabularx}

\begin{nota}
Para projetos de legislação municipal em português como este, a Maritaca tende a produzir
respostas mais naturais e precisas no idioma.
\end{nota}

% ===========================================================================
\section{Instalação após as alterações}

Se você já tinha o projeto instalado, atualize as dependências para incluir o SDK
\texttt{openai} (necessário para a Maritaca):

\begin{lstlisting}
pip install -r bh_news_chatbot/requirements.txt
\end{lstlisting}

Ou instale apenas o novo pacote:

\begin{lstlisting}
pip install "openai>=1.0.0"
\end{lstlisting}

% ===========================================================================
\section{Perguntas frequentes}

\subsection{Os testes de erros passam mesmo sem SDK instalada. Isso está certo?}

Sim. Testes como ``chave inválida $\to$ complete() falha'' verificam apenas que
\emph{alguma exceção é lançada}. Sem o SDK instalado, o Python lança
\texttt{ModuleNotFoundError} antes mesmo de chamar a API --- o teste ainda passa porque o
comportamento esperado (falha) ocorreu. Com o SDK instalado e chave inválida, seria lançado
um \texttt{AuthenticationError} da API.

\subsection{Por que os testes funcionais ficam como \texttt{-} (pulado)?}

Porque a chave de API não está configurada. Configure a variável de ambiente e rode
novamente:

\begin{lstlisting}
export ANTHROPIC_API_KEY=sk-ant-...
python bh_news_chatbot/test_api.py --provider anthropic
\end{lstlisting}

\subsection{Posso usar mais de um provedor ao mesmo tempo?}

Não simultaneamente --- você escolhe um por execução com \texttt{-{}-provider}. Mas pode
ter as três chaves configuradas e alternar conforme precisar.

\subsection{Como saber qual provedor está sendo usado no chat?}

O banner inicial do chatbot exibe o nome do provedor ativo:

\begin{lstlisting}
Decretos carregados (Maritaca (sabia-4)):
  * Area X: 10 documento(s)
  ...
\end{lstlisting}

\subsection{O teste ``modelo inexistente'' fica como \texttt{-} mesmo com chave válida. Por quê?}

Verifique se a chave está configurada com \texttt{export} e não apenas como assignment
simples:

\begin{lstlisting}
export ANTHROPIC_API_KEY=sk-ant-...   # correto
ANTHROPIC_API_KEY=sk-ant-...          # errado -- nao fica no ambiente do subprocesso
\end{lstlisting}

\subsection{O que fazer se \texttt{complete()} falhar com erro de autenticação?}

\begin{enumerate}
  \item Verifique se a chave está correta e sem espaços extras
  \item Verifique se a chave não expirou ou foi revogada
  \item Para Anthropic: confirme créditos em \href{https://console.anthropic.com}{console.anthropic.com}
  \item Para Maritaca: confirme o saldo em \href{https://plataforma.maritaca.ai}{plataforma.maritaca.ai}
  \item Para Gemini: confirme que a API está habilitada no \href{https://aistudio.google.com}{Google AI Studio}
\end{enumerate}

\subsection{Como usar o modelo rápido para reduzir custos?}

O modelo rápido é usado automaticamente na categorização interna. Para o chat, passe
\texttt{-{}-model} explicitamente:

\begin{lstlisting}
# Anthropic -- modelo barato
python bh_news_chatbot/main.py --provider anthropic --model claude-haiku-4-5-20251001

# Maritaca -- modelo barato
python bh_news_chatbot/main.py --provider maritaca --model sabiazinho-4

# Gemini -- flash ja e o padrao
python bh_news_chatbot/main.py --provider gemini --model gemini-2.0-flash
\end{lstlisting}

\subsection{O token do GitHub é necessário para usar o chatbot?}

Não. O token do GitHub é usado apenas para operações no repositório (push, pull, consultas
à API). O chatbot precisa apenas das chaves de LLM: \texttt{ANTHROPIC\_API\_KEY},
\texttt{GEMINI\_API\_KEY} ou \texttt{MARITACA\_KEY}.

\begin{nota}
\textbf{Importante:} nunca compartilhe tokens ou chaves de API em conversas, issues ou
commits. Se isso acontecer, revogue imediatamente o token no painel do serviço correspondente.
\end{nota}

% ===========================================================================
\section{Referência rápida de comandos}

\begin{lstlisting}
# Verificar APIs antes de usar
python bh_news_chatbot/test_api.py

# Inspecionar o cache de artigos
python bh_news_chatbot/debug.py

# Iniciar o chatbot (Anthropic)
python bh_news_chatbot/main.py

# Iniciar o chatbot (Gemini)
python bh_news_chatbot/main.py --provider gemini

# Iniciar o chatbot (Maritaca)
python bh_news_chatbot/main.py --provider maritaca

# Forcar re-download dos artigos
python bh_news_chatbot/main.py --refresh

# Categorizar artigos e ver relatorio
python bh_news_chatbot/main.py --only-categorize

# Limpar cache manualmente
rm bh_news_chatbot/cache/news_cache.json
rm bh_news_chatbot/cache/categories_cache.json
\end{lstlisting}

\end{document}
